%=====================================================================
\section{Schichtmodelle und Scheibenzerlegung}\label{sec:layers}
%=====================================================================

\subsection{Schichttypen}
\begin{center}
\small
\begin{tabular}{p{3.2cm}p{4.7cm}p{6.2cm}}
\toprule
Typ & Beschreibung & $\varepsilon$ bei Höhe $y$ / Normalenfeld \\
\midrule
\verb|Homogeneous| & Dicke $h$, Material & konstant; analytische Moden, keine Integration \\
\verb|Lamellar1D| & Intervalle $[x_j,x_{j+1}]$ mit Materialien, Dicke $h$ & \eqref{eq:interval}; $N_x^2=1$ \\
\verb|Profile1D| & $y=g(x)$ (Funktion oder Stützstellen), Material unten/oben, Tiefe $a$ &
 Schnittpunkte $g(x)=y$ per Nullstellensuche $\Rightarrow$ Intervalle; $\mathbf N$ aus \eqref{eq:N1d} \\
\verb|Polygon2D| & $y$-invarianter Querschnitt aus Polygonen/Rechtecken/Kreisen, Dicke $h$ &
 \eqref{eq:polygon}; $\mathbf N$ nach Algorithmus NF \\
\verb|Profile2D| & $y=g(x,z)$, Material unten/oben, Tiefe $a$ & Abtastung von $\{g>y\}$ + Polygonisierung
 (Marching Squares) oder feines FFT-Raster; $\mathbf N$ aus \eqref{eq:N2d} \\
\verb|Function| & benutzerdefiniertes $\overline{\overline\varepsilon}(x,y,z)$ & Abtastung + FFT (nur Tests) \\
\bottomrule
\end{tabular}
\end{center}
Beschichtete Profile (mehrere parallele Profile $g(x)+t_k$) werden als ein Schichtobjekt mit mehreren
Grenzflächen behandelt; $\mathbf N$ ist dann für alle Grenzflächen identisch (Buch IV.2.5).

\subsection{Scheibenzerlegung}
Jede modulierte Schicht wird in $K$ Scheiben $[y_{k-1},y_k]$ zerlegt (gleichmäßig oder nach Krümmung
von $g$). Kriterium gegen Kontamination (Buch III.6, mit Sicherheitsabstand):
\begin{equation}
  \rho_{\max}\,(y_k-y_{k-1})\le\ln(10^{8})\approx18.4,\qquad
  \rho_{\max}\approx\max_{(n,m)}\big|\sqrt{\alpha_n^2+\gamma_m^2-\varepsilon_{\min}}\big|,
  \label{eq:slicecrit}
\end{equation}
mit $\varepsilon_{\min}$ = Material mit dem kleinsten $\operatorname{Re}\varepsilon$ (Metalle: großes $|\rho|$).
\textbf{Achtung (Befund M3):} Die abgeschnittene FFF-Matrix kann Eigenwerte weit über dieser
Materialschätzung haben (Gold-Sinus: $|\lambda|\approx170$ statt $\approx6$). Der RK-Löser verwendet daher
$\rho_{\max}=\max(\text{Schätzung},\ \max_k|\lambda(M(y_k))|)$ an fünf Höhen $y_k$ und halbiert eine Scheibe
zusätzlich, solange $\max|T_F|>10^{8}$ (a-posteriori-Kontrolle).
RCWA: zusätzlich Treppenfehler $\Rightarrow$ $K$ per Konvergenztest. RK: Scheiben nur zur Stabilisierung,
die Genauigkeit regelt die Schrittweite.

\subsection{Permittivität bei beliebiger Höhe}
Der RK-Löser braucht $\tp{\varepsilon}(y)$, $\tp{1/\varepsilon}(y)$ bzw.\ $\tp{C_\varepsilon}(y)$ an den
Stufenknoten $y+c_ih$. \verb|Profile1D|: Nullstellen von $g(x)-y$ im Intervall $[0,d_x)$ (Abtasten auf
feinem Gitter, Vorzeichenwechsel, dann Brent/Newton), daraus Intervalle und exakte Koeffizienten
\eqref{eq:interval}. \verb|Profile2D|: Höhenlinie $g=y$ per Marching Squares als Polygon, dann
\eqref{eq:polygon}; Fallback FFT mit Oversampling. Cache: Koeffizienten je (Schicht, $y$) speichern, da
bei Feldrekonstruktion dieselben Knoten erneut gebraucht werden.

%=====================================================================
\section{Löser A: implizites Runge--Kutta (Differentialmethode)}\label{sec:rk}
%=====================================================================

\subsection{Verfahren}
Lineares System \eqref{eq:ode}. Ein $s$-stufiges RK-Verfahren mit Butcher-Tableau $(\mathsf A,\mathbf b,\mathbf c)$:
\begin{equation}
  K_i=M_i\Big(F+h\sum_{j=1}^{s}a_{ij}K_j\Big),\quad M_i=M(y+c_ih),\qquad
  F(y+h)=F+h\sum_{i=1}^{s}b_iK_i .
  \label{eq:rk}
\end{equation}
Für lineare Systeme ergibt sich ein \textbf{Schrittpropagator} (unabhängig von $F$):
\begin{equation}
  \big(\mathbb 1_{sn}-h\,\mathfrak A\big)X=\mathfrak M,\qquad
  \mathfrak A=\big(a_{ij}M_i\big)_{i,j=1}^{s},\quad \mathfrak M=\begin{pmatrix}M_1\\\vdots\\M_s\end{pmatrix},\qquad
  P_h=\mathbb 1_n+h\sum_{i=1}^{s}b_iX_i ,
  \label{eq:rkprop}
\end{equation}
mit $n=4n_f$ (bzw.\ $2n_f$ in 1D-in-plane), $X=(X_1;\dots;X_s)$, $X_i\in\mathbb C^{n\times n}$. Kosten pro
Schritt: eine LU-Zerlegung der Größe $sn$ und eine Lösung mit $n$ rechten Seiten.
Die Scheiben-Transfermatrix im Feldraum ist $T_F=\prod_kP_{h_k}$ (in Reihenfolge von unten nach oben).

\paragraph{Gauß--Legendre (Standard), $s=2$, Ordnung 4:}
\[
  \mathbf c=\Big(\tfrac12-\tfrac{\sqrt3}6,\ \tfrac12+\tfrac{\sqrt3}6\Big),\quad
  \mathsf A=\begin{pmatrix}\tfrac14&\tfrac14-\tfrac{\sqrt3}6\\\tfrac14+\tfrac{\sqrt3}6&\tfrac14\end{pmatrix},\quad
  \mathbf b=\big(\tfrac12,\tfrac12\big).
\]
\paragraph{Gauß--Legendre, $s=3$, Ordnung 6:}
\[
  \mathbf c=\Big(\tfrac12-\tfrac{\sqrt{15}}{10},\tfrac12,\tfrac12+\tfrac{\sqrt{15}}{10}\Big),\quad
  \mathsf A=\begin{pmatrix}\tfrac5{36}&\tfrac29-\tfrac{\sqrt{15}}{15}&\tfrac5{36}-\tfrac{\sqrt{15}}{30}\\
  \tfrac5{36}+\tfrac{\sqrt{15}}{24}&\tfrac29&\tfrac5{36}-\tfrac{\sqrt{15}}{24}\\
  \tfrac5{36}+\tfrac{\sqrt{15}}{30}&\tfrac29+\tfrac{\sqrt{15}}{15}&\tfrac5{36}\end{pmatrix},\quad
  \mathbf b=\big(\tfrac5{18},\tfrac49,\tfrac5{18}\big).
\]
\paragraph{Radau IIA, $s=3$, Ordnung 5 (optional, L-stabil):}
\[
  \mathbf c=\Big(\tfrac{4-\sqrt6}{10},\tfrac{4+\sqrt6}{10},1\Big),\quad
  \mathsf A=\begin{pmatrix}\tfrac{88-7\sqrt6}{360}&\tfrac{296-169\sqrt6}{1800}&\tfrac{-2+3\sqrt6}{225}\\
  \tfrac{296+169\sqrt6}{1800}&\tfrac{88+7\sqrt6}{360}&\tfrac{-2-3\sqrt6}{225}\\
  \tfrac{16-\sqrt6}{36}&\tfrac{16+\sqrt6}{36}&\tfrac19\end{pmatrix},\quad \mathbf b=\text{letzte Zeile von }\mathsf A .
\]
\textbf{Tests:} Ordnungsbedingungen ($\sum b_i=1$, $\sum b_ic_i=\tfrac12$, $\sum b_ic_i^2=\tfrac13$, \dots\ bis zur
Ordnung), $\mathsf A\mathbf 1=\mathbf c$; für konstantes $M$ ist $P_h$ bei Gauß--Legendre der
$(s,s)$-Padé-Approximant von $e^{hM}$ $\Rightarrow$ Vergleich mit \verb|expm|, Fehler $O(h^{2s+1})$ je Schritt.

\subsection{Stabilität und Schrittweite}
Gauß--Legendre ist A-stabil, $|R(z)|=1$ auf der imaginären Achse (propagierende Moden ohne
künstliche Dämpfung). Evaneszente Moden ($\lambda$ reell, groß) werden nur für $h|\lambda|\lesssim1$ genau
wiedergegeben; zu große Schritte verfälschen sie (bei Gauß nicht gedämpft, $|R(\pm\infty)|=1$).
Daher:
\begin{equation}
  h\le c_s/\rho_{\max},\qquad c_2\approx0.5,\ c_3\approx1.0,
  \label{eq:stepcrit}
\end{equation}
plus Fehlersteuerung per Schrittverdopplung: $\epsilon=\|P_hP_h-P_{2h}\|_F/\|P_{2h}\|_F$ mit
$P_h$ über $[y,y+h]$, $[y+h,y+2h]$; Ziel $\epsilon\le\texttt{rk\_tol}$ (Standard $10^{-8}$),
Schrittweitenanpassung $h\leftarrow h\,(\texttt{rk\_tol}/\epsilon)^{1/(2s+1)}\cdot0.9$.
Für $y$-unabhängiges $M$ (Lamellen, homogene Bereiche) genügt eine Padé-Auswertung mit Skalieren
und Quadrieren: $P=(P_{h/2^k})^{2^k}$.

\subsection{Scheiben-S-Matrix aus dem RK-Ergebnis}
\begin{equation}
  T^{(\text{Scheibe})}=\Psi_{\text{ref}}^{-1}\,T_F\,\Psi_{\text{ref}},\qquad S^{(\text{Scheibe})}\ \text{aus \eqref{eq:T2S}}
  \label{eq:rkslice}
\end{equation}
(Buch Anh.~IV (A.IV.14'); dort Blockgröße für 2D falsch angegeben: richtig $n_h=2n_f$).
Die Scheibendicke folgt \eqref{eq:slicecrit}, damit $T_{11}$ gut konditioniert ist; Diagnose:
Konditionszahl von $T_{11}$ schätzen und bei $>10^{10}$ die Scheibe halbieren.

\subsection{Dichte Ausgabe für das Nahfeld}
Gauß--Kollokation liefert ein Polynom vom Grad $s$ je Schritt:
\begin{equation}
  F(y+\vartheta h)=F(y)+h\sum_{j=1}^{s}\Big(\int_0^{\vartheta}\ell_j(\tau)\,d\tau\Big)K_j,\qquad 0\le\vartheta\le1,
  \label{eq:dense}
\end{equation}
mit den Lagrange-Polynomen $\ell_j$ zu den Knoten $\mathbf c$. Für das Nahfeld wird nur \emph{ein}
Zustandsvektor (die tatsächliche Lösung) von der Unterkante jeder Scheibe aus neu integriert
(dieselbe Schrittfolge, gespeichert in \verb|RKSliceRecord|).

\subsection{Rechenaufwand (Richtwerte)}
Pro Schritt $\approx\frac83(sn)^3$ reelle Flops (LU) $+8(sn)^2n$ (Lösen). 2D, $N_x=N_z=10$:
$n=1764$, $s=2$: $\approx4.7\cdot10^{11}+2\cdot10^{11}$ Flops. Der RK-Löser ist daher für 2D teuer und
hauptsächlich für Profile $g(x,z)$ und Referenzrechnungen gedacht; für $y$-invariante 2D-Schichten
ist RCWA vorzuziehen.

%=====================================================================
\section{Löser B: RCWA (Eigenmoden je Scheibe)}\label{sec:rcwa}
%=====================================================================

\subsection{Eigenwertproblem}
In einer Scheibe mit $y$-unabhängigem $M$ (Treppenapproximation: $\varepsilon$ in Scheibenmitte):
$M=W\Lambda W^{-1}$, $F(y)=\sum_kc_kw_ke^{\lambda_k(y-y_{\text{ref},k})}$.
\paragraph{Reduziertes Problem (isotrop/lamellar, $M_{11}=M_{12}=M_{21}=M_{22}=0=M_{33}=M_{34}=M_{43}=M_{44}$):}
Mit $\mathcal P=\begin{pmatrix}M_{13}&M_{14}\\M_{23}&M_{24}\end{pmatrix}$, $\mathcal R=\begin{pmatrix}M_{31}&M_{32}\\M_{41}&M_{42}\end{pmatrix}$:
\begin{equation}
  M=\begin{pmatrix}0&\mathcal P\\\mathcal R&0\end{pmatrix},\quad
  \mathcal P\mathcal R\,u=\mu\,u,\quad \kappa=\sqrt{-\mu}\ (\text{Zweig \eqref{eq:beta_branch}}),\quad
  w_\pm=\begin{pmatrix}u\\\mathcal Ru/\lambda_\pm\end{pmatrix},\ \lambda_\pm=\pm i\kappa .
  \label{eq:rcwa_reduced}
\end{equation}
(Prüfung: $Mw_\pm=(\mathcal P\mathcal Ru/\lambda_\pm;\mathcal Ru)=\lambda_\pm w_\pm$ wegen $\mu=\lambda_\pm^2$.)
\enquote{Auf}: $\lambda_+=i\kappa$ ($\operatorname{Re}\lambda_+\le0$), \enquote{ab}: $\lambda_-=-i\kappa$.
Für $\kappa\approx0$ (Grenzfall) $w=(u;\mathcal Ru/\lambda)$ schlecht definiert $\Rightarrow$ kleine
Regularisierung wie bei \eqref{eq:beta_branch}.
\paragraph{Allgemein (anisotrop, Profile mit $Q_{yx}\neq0$):} volle Eigenzerlegung der $4n_f$-Matrix.
Klassifikation (robuste Form von \eqref{eq:classify}): Bewertung $s_k=\operatorname{Re}\lambda_k$, für
$|\operatorname{Re}\lambda_k|\le\tau$ stattdessen $\mp\tau/2$ nach dem Vorzeichen von $S_y$; sortieren und die untere
Hälfte als aufwärts nehmen. (Die reine Schwellenwert-Variante scheiterte in metallischen Scheiben.)
\paragraph{TE/TM-Entkopplung:} Im 1D-in-plane-Fall verschwinden alle Blöcke zwischen $(E_x,H_z)$ und
$(E_z,H_x)$ exakt. RK und RCWA erkennen das automatisch und rechnen beide Hälften getrennt
(Laufzeitgewinn RK $\approx6$, RCWA $\approx2{,}5$).

\subsection{Stabile Scheiben-S-Matrix}
Moden auf ($W_u,\Lambda_u$, Bezug Unterkante $y_u$) und ab ($W_d,\Lambda_d$, Bezug Oberkante $y_o$):
\begin{equation}
  F(y)=W_ue^{\Lambda_u(y-y_u)}\mathbf c_u+W_de^{\Lambda_d(y-y_o)}\mathbf c_d,\qquad
  X_u=e^{\Lambda_uh},\ X_d=e^{-\Lambda_dh}\quad(|X|\le1).
  \label{eq:rcwa_field}
\end{equation}
Mit $\Psi_{\text{ref}}^{-1}W_u=\begin{pmatrix}A_u\\B_u\end{pmatrix}$, $\Psi_{\text{ref}}^{-1}W_d=\begin{pmatrix}A_d\\B_d\end{pmatrix}$
(obere Hälfte: ab-Amplitude, untere: auf-Amplitude) gilt an Unter- und Oberkante
\begin{equation}
  \begin{aligned}
  \mathbf a_u&=A_u\mathbf c_u+A_dX_d\mathbf c_d, & \mathbf b_u&=B_u\mathbf c_u+B_dX_d\mathbf c_d,\\
  \mathbf a_o&=A_uX_u\mathbf c_u+A_d\mathbf c_d, & \mathbf b_o&=B_uX_u\mathbf c_u+B_d\mathbf c_d .
  \end{aligned}
  \label{eq:rcwa_bc}
\end{equation}
Die Eingänge $(\mathbf b_u,\mathbf a_o)$ bestimmen
\begin{equation}
  \underbrace{\begin{pmatrix}B_u&B_dX_d\\A_uX_u&A_d\end{pmatrix}}_{=:G}\begin{pmatrix}\mathbf c_u\\\mathbf c_d\end{pmatrix}=
  \begin{pmatrix}\mathbf b_u\\\mathbf a_o\end{pmatrix},\qquad
  \begin{pmatrix}\mathbf b_o\\\mathbf a_u\end{pmatrix}=
  \underbrace{\begin{pmatrix}B_uX_u&B_d\\A_u&A_dX_d\end{pmatrix}}_{=:H}\begin{pmatrix}\mathbf c_u\\\mathbf c_d\end{pmatrix},
  \qquad S=HG^{-1}.
  \label{eq:rcwa_S}
\end{equation}
Alle Einträge sind beschränkt; keine wachsenden Exponentiale. $G$ als LU speichern
(\verb|RCWASliceRecord|) für die Feldrekonstruktion. Homogene Zwischenschichten: dieselben Formeln mit
analytischen Moden ($W=\Psi_X$, $\Lambda=\operatorname{diag}(\mp i\beta)$).

\subsection{Querprüfung RK $\leftrightarrow$ RCWA}
Für $y$-unabhängiges $M$ (Lamellen) müssen beide Löser bis auf die RK-Toleranz identische Effizienzen
liefern; für glatte Profile müssen RCWA mit $K\to\infty$ Scheiben und RK gegen denselben Wert
konvergieren -- \emph{außer} im TM-Fall bei Metallen mit Li-Lamellen-Faktorisierung (Buch VI.5:
Treppenapproximation konvergiert dort extrem langsam). Deshalb verwendet RCWA auf Profilschichten
\textbf{dieselbe FFF-Faktorisierung mit dem Profil-Normalenfeld} ($\varepsilon$ in Scheibenmitte, $\mathbf N$ des
glatten Profils), nicht die Lamellenformel. Kapitelprüfung: Al-Sinusgitter, FFF mit 60/150/300 Scheiben
$\to0{,}853/0{,}865/0{,}869$, Referenz $0{,}8696$ (C-Methode).

%=====================================================================
\section{Nahfeld}\label{sec:field}
%=====================================================================

\subsection{Ablauf}
\begin{enumerate}[nosep]
  \item Gesamtlösung \eqref{eq:solution}, interne Amplituden $(\mathbf a_q,\mathbf b_q)$ an allen Grenzflächen \eqref{eq:internal}.
  \item Fourierkoeffizienten $F(y)=([E_x],[E_z],[H_x],[H_z])$ an den gewünschten Höhen:
    \begin{itemize}[nosep]
      \item homogene Schicht/Superstrat/Substrat: Moden mit Phasen $e^{\lambda(y-y_{\text{ref}})}$;
      \item RCWA-Scheibe: $(\mathbf c_u;\mathbf c_d)=G^{-1}(\mathbf b_u;\mathbf a_o)$, dann \eqref{eq:rcwa_field};
      \item RK-Scheibe: $F(y_u)=\Psi_{\text{ref}}(\mathbf a_u;\mathbf b_u)$, erneute Integration mit dichter Ausgabe \eqref{eq:dense}.
    \end{itemize}
  \item $[E_y]$, $[H_y]$ aus \eqref{eq:Ey}, \eqref{eq:Hy} mit den $Q$-Matrizen der Höhe $y$.
  \item Synthese im Ortsraum: $u(x,y,z)=\sum_ru_r(y)e^{i(\alpha_nx+\gamma_mz)}$, effizient per inverser FFT
        auf einem Gitter (Zero-Padding), danach Multiplikation mit $e^{i(\alpha_0x+\gamma_0z)}$.
\end{enumerate}

\subsection{Umsetzung (M6)}
\begin{itemize}[nosep]
  \item Ein Lauf mit Aufzeichnung speichert jedes Element des Stapels (Grenzfläche zum Substrat, modale Scheibe
        mit $W$, $\Lambda$, $G$, RK-Scheibe mit den Transfermatrizen $T(y_0\to y_j)$ an genau den angeforderten
        Höhen $y_j$). Die RK-Integration hält dafür an den Abtasthöhen an; dichte Ausgabe \eqref{eq:dense}
        ist damit nicht nötig.
  \item Interne Amplituden aus \eqref{eq:internal} mit kumulierten $S_L$, $S_U$ je Grenzfläche.
  \item Schnittstelle \verb|Solver::fieldCoefficients(stack, wl, y)| liefert $F$, $[E_y]$, $[H_y]$ je Höhe;
        \verb|Solver::nearField| erzeugt die Schnitte. Gemeinsame Schichtmodelle (\verb|detail::LayerModel|:
        $Q(y)$, $\overline{\overline\varepsilon}(x,y)$, $\mathbf N(x)$, $\tp{\mathbf N\mathbf N^T}$) sorgen dafür, dass
        Lösung und Rekonstruktion dieselben Daten verwenden (RCWA-Treppe: $Q$ und $\varepsilon$ der Scheibenmitte).
\end{itemize}

\subsection{FFF-konsistente Rekonstruktion (empfohlen in modulierten Schichten)}
Die direkte Synthese von $E_x,E_y$ aus abgeschnittenen Reihen zeigt Gibbs-Oszillationen an
Materialgrenzen. Besser: stetige Größen synthetisieren und punktweise zurückrechnen:
\begin{equation}
  [\mathbf F_\varepsilon]=\tp{C_\varepsilon}^{-1}[\mathbf E]\ \ \text{(Variante A)}\ \Rightarrow\ \mathbf F_\varepsilon(x,z)\ \text{(stetig)}
  \ \Rightarrow\ \mathbf E(x,z)=C_\varepsilon(x,z)\,\mathbf F_\varepsilon(x,z),
  \label{eq:recon}
\end{equation}
umgesetzt (für isotrope und anisotrope Schichten gleichermaßen) über die Projektorform \eqref{eq:projector}:
$[\mathbf D]=Q[\mathbf E]$, $[\mathbf d_n]=\tp{P_N}[\mathbf D]$, $[\mathbf e_t]=[\mathbf E]-\tp{P_N}[\mathbf E]$ synthetisieren
(stetig) und punktweise $\mathbf E(x)=H_1(x)\,\mathbf e_t(x)+\xi(x)^{-1}\mathbf d_n(x)$; isotrop ist das
$\mathbf E=\mathbf e_t+\mathbf d_n/\varepsilon$. Befund (V12): an einer Lamellenwand springt $D_x=\varepsilon E_x$ bei
direkter Synthese um 270\,\%, FFF-konsistent um 3\,\% (Gitterabstand 2\,nm). $\mathbf H$ ist überall stetig ($\mu=1$) und wird direkt synthetisiert.
Beide Varianten (direkt / FFF-konsistent) sind auswählbar (\verb|reconstruction: direct \| fff|).

\subsection{Schnitte und Ausgabe}
\begin{center}
\small
\begin{tabular}{lll}
\toprule
& horizontaler Schnitt ($y=y_0$) & vertikaler Schnitt \\
\midrule
1D & Linie $u(x,y_0)$, $x\in[x_a,x_b]$ & Karte $u(x,y)$, $z=0$ \\
2D & Karte $u(x,y_0,z)$ & Karte $u(x,y,z_0)$ (optional zusätzlich $u(x_0,y,z)$) \\
\bottomrule
\end{tabular}
\end{center}
Größen: alle sechs komplexen Komponenten, $|\mathbf E|^2$, $|\mathbf H|^2$, $\operatorname{Re}(\mathbf Ee^{-i\omega t})$ für
gegebene Phase, Poynting-Vektor $\tfrac12\operatorname{Re}(\mathbf E\times\mathbf H^*)$, Materialindex-Karte (für Konturen).
Ausgabeformat: je Schnitt ein Verzeichnis (\verb|nearfield_horizontal|, \verb|nearfield_vertical|) mit
\verb|Ex|\dots\verb|Hz.npy| (complex128), \verb|absE2|, \verb|absH2|, \verb|Sy|, \verb|material.npy| (float64),
C-Ordnung, Form $(n_y,n_x)$ bzw.\ $(n_x)$, plus \verb|meta.json|; $\mathbf H$ als $Z_0\mathbf H$, $|\mathbf E_{\text{inc}}|=1$ (Achsen in µm, Wellenlänge,
Einfall, Komponentenliste). Optional HDF5 (\verb|FG_ENABLE_HDF5|). Das Skript
\verb|scripts/plot_nearfield.py| erzeugt PNG/PDF.
